\section{Sparse standard geometry in free quadrilateral coordinates}
\label{sec:window-geometry}

The propagated potentials of Section~\ref{sec:window-modes} already have
only $d$ parameters. The old window constructor expands each form into
$k$ coefficients and every standard constraint into $p+k$ coefficients,
where $p$ is the number of retained triangle classes. Low-dimensional
producers then multiply those expansions, including their zeros. This
continuation keeps the actual coefficients sparse and materializes dense
rows only for rank computations that request them.

\subsection{The unchanged complete matching equations}

Let $F_a$ be the free Q coordinates of the canonical updated basis $K'$,
so $q_{F_a}=z_a$. Every dependent Q coordinate obeys
\[
                 q_j-\sum_{a=1}^d K'_{a,j}q_{F_a}=0.
\]
Each such equation has at most $d+1$ nonzero coefficients. A class potential
$\phi_c(z)$ has at most $d$ coefficients, all in the free physical slots.
For a vertex-group anchor $b$, the standard equation is
\[
 x_c-x_b-\sum_{a=1}^d(\phi_{c,a}-\phi_{b,a})q_{F_a}=0,
\]
with at most $d+2$ nonzeros. Together these are the same complete standard
matching space as the preceding projected kernel, including its link-offset
factors. The representation change neither drops constraints nor replaces
standard rays with Q rays.

\subsection{Exact forms with two access paths}

The native \code{SparseLinearForm} stores the width and the nonzero
coefficient dictionary. Its sparse dot product visits only stored entries.
Its chart projection visits the same entries for each chart coordinate.
Indexing returns the exact coefficient, including omitted zeros; iteration
returns the full dense sequence. Coordinate-face restrictions can therefore
index the sparse row directly, while the existing rational rank routines
can materialize the same dense row by iteration.

\begin{proposition}[Representation preserves all kernel queries]
Sparse window constraints and potentials define exactly the previous
projected kernel. Their sparse dot and chart projection values equal dense
arithmetic; every coordinate-face restriction and dense rank calculation
receives exactly the previous coefficient values.
\end{proposition}
\begin{proof}
Only exact zero coefficients are omitted. Summation over nonzero entries
is equal to summation over the full row, for every vector or chart form.
Indexing and iteration restore zero at each omitted slot and retain every
other value. Thus both matching equations and restricted matrices are
identical. The basis, normalized potentials, canonical lifts, extrema and
rank/nullspace tests consequently remain unchanged.
\end{proof}

The projected lift evaluates its Q constraints and potentials sparsely,
then checks integrality and validates the full coordinate vector against
the actual source as before. Envelope line construction uses sparse dot
products. Planar chart construction uses sparse form projection. Other
kernel constructors still use their existing dense tuples. Their fallback
arithmetic is retained. The generic arrangement/support paths see the
complete restored rows when needed; no low-dimensional promise is applied
to a higher-nullity source.

\subsection{Storage and arithmetic accounting}

There are $k-d$ dependent-Q rows and $p-g$ triangle-difference rows, with
$g$ retained vertex groups. The stored standard coefficient count is at most
\[
          (k-d)(d+1)+(p-g)(d+2).
\]
Stored potentials use at most $pd$ entries. Thus standard coefficient and
potential storage is $O((k+p)(d+1))$, instead of expanded row storage with
quadratic dependence on $k+p$. Widths and basis storage remain separately
charged. Source class identification still uses its existing exact
form-group dictionary; its hash costs are not replaced by a claimed
deterministic linear-time bound.

After the basis and propagated forms are available, filling the sparse
coefficient records requires $O((k+p)(d+1))$ scalar operations. Checking a
Q vector costs $O(k(d+1))$ sparse coefficient operations. Canonical potential
evaluation costs $O(pd)$, followed by the existing minima, physical-corner
reconstruction and independent source validation. A one-dimensional
direction therefore avoids the old dense $k^2$ cycle-row dot products.
Envelope and planar projections likewise avoid multiplying zero physical
Q slots; their fan, clipping and actual-ray costs remain charged.

Generic Gaussian elimination still pays for any dense matrix it constructs.
Exact coefficient bits, indexing/hash costs, output and certificate replay
also remain charged. With logarithmic matching nullity this reduces
storage and a polynomial factor of the supplied-window computation. It
supplies neither the missing global window family nor a general
quasi-polynomial recognition theorem.

\subsection{Source validation and measured behavior}
All 1,411 maintained tests pass in 262.025 seconds; the thirty-eight focused
tests pass in 23.105 seconds. Genuine nullity-one, two and three sectors
produce the complete expected source rays while dense-form iteration is
disabled. Direct projected lifts and coordinate-face extremality tests
also work under that restriction. A genuine nullity-four sector retains
both complete generic methods. A larger sixteen-layer source verifies
the linear stored-coefficient bounds and exact lifts scaled by $2^{90}+1$.

The 85-source audit exactly preserves every disabled output, every enabled
verdict, every recorded guard count and all 31 native positives. Seven
ordinary window discs receive fresh Regina confirmation. The twenty-eight
capped cases and twenty-six completed local misses remain unchanged.
All 602 captured runtime pins remain fixed. The 214.396-second audit is
not a paired speedup estimate against its predecessor. Its result establishes
preservation of the configured source behavior, not a global knot-search
completeness theorem.

Native timings freeze seven preceding-release modules at
\code{9abd16fe1}: sector dispatch, envelope and planar production, window
basis construction, residual queries, seed queries and recognition.
Local imports are redirected only to the matching frozen dependencies,
including the generic dispatcher's lazy low-dimensional imports. Five
serial interleaved four-arm rounds complete 200 measured calls and forty
warmups. The timings include source and cache preparation, complete
configured queries, independent positive source replay and serialized output.
Every arm completes its configured work. Statuses and tested positive
certificate digests agree exactly.

\begin{samepage}
\noindent\textbf{Complete native queries with radius-two windows.}
\begin{center}
\small
\begin{tabular}{lrrrrr}
\toprule
Input & previous ms & current ms & old/new & old A/A & new A/A\\
\midrule
empty circle & 5.923 & 6.068 & 0.987 & 0.994 & 1.016\\
optimized positive & 63.597 & 60.929 & 1.012 & 1.003 & 0.991\\
genus one miss & 237.021 & 229.417 & 1.024 & 0.995 & 0.976\\
trefoil & 414.162 & 393.843 & 1.044 & 1.017 & 1.013\\
figure eight & 722.362 & 679.206 & 1.064 & 0.986 & 1.000\\
\bottomrule
\end{tabular}
\end{center}

\end{samepage}

The native genus-one, trefoil and figure-eight queries have paired gains
of about 1.024, 1.044 and 1.064. Early controls that return before sparse
window geometry remain close to parity. These are improvements over the
preceding cached-mode implementation, not over disabling the window stage.

\begin{samepage}
\noindent\textbf{Complete recognition with radius-two windows enabled.}
\begin{center}
\small
\begin{tabular}{lrrrrr}
\toprule
Input & previous ms & current ms & old/new & old A/A & new A/A\\
\midrule
empty circle & 0.018 & 0.018 & 1.000 & 0.977 & 0.967\\
optimized positive & 51.249 & 51.234 & 0.996 & 1.006 & 1.007\\
genus one miss & 244.378 & 230.930 & 1.039 & 1.004 & 0.986\\
trefoil & 435.224 & 407.538 & 1.063 & 1.008 & 0.975\\
figure eight & 754.915 & 716.218 & 1.052 & 0.996 & 0.960\\
\bottomrule
\end{tabular}
\end{center}

\end{samepage}

Complete recognition has paired gains of 1.039, 1.063 and 1.052 on the
same three cases. The first figure-eight new-arm A/A control is 0.960,
so a separate fifteen-round comparison repeats the two knot cases with
full recognition and serialized output. It completes 120 measured calls
and eight warmups. The trefoil gain is 1.046, with A/A controls 0.985 and
0.994; the figure-eight gain is 1.052, with controls 0.996 and 1.003.
Both repeat rows include the complete fallback after the native miss.
The longer records support modest measured gains, rather than a uniform
wall-clock theorem. The microsecond empty-recognition control is not used
to infer a geometry speedup. Radius zero remains much faster on the small
controls, and the window stage stays off by default.

\begin{sloppypar}
Runtime release \code{f3772d8a3}, tests, the source audit, final timing
run and longer repeat are retained under \code{data/window-geometry-*}.
From \code{fast/}, use \code{normal\_orbit\_research.window\_geometry} with
\code{audit --fresh-regina} or \code{benchmark --rounds 5}.
Publication checks compare 602 frozen/current runtime pins, 599 preceding
release evidence pins, every timing outcome and all 31 saved positive proofs.
The proof replay disables sparse dot/chart evaluation, cached modes,
Euler compilation, prepared observers, matching updates, residual planning,
coherent extraction, ray enumeration and orbit producers. Independent source
geometry remains the authority for positive results.
\end{sloppypar}

